\documentclass[11pt, a4paper]{article}

% ── Packages ──────────────────────────────────────────────────────────
\usepackage[utf8]{inputenc}
\usepackage[T1]{fontenc}
\usepackage{lmodern}
\usepackage[margin=1in]{geometry}
\usepackage{amsmath, amssymb, amsthm}
\usepackage{booktabs}
\usepackage{longtable}
\usepackage{graphicx}
\usepackage{hyperref}
\usepackage{xcolor}
\usepackage{listings}
\usepackage{tcolorbox}
\usepackage{polya}

\hypersetup{
  colorlinks=true,
  linkcolor=polyablue,
  urlcolor=polyablue,
  citecolor=polyablue,
}

% ── Code listing style ────────────────────────────────────────────────
\lstdefinestyle{polya-python}{
  language=Python,
  basicstyle=\ttfamily\small,
  keywordstyle=\color{polyablue}\bfseries,
  commentstyle=\color{polyagray}\itshape,
  stringstyle=\color{polyagreen},
  numbers=left,
  numberstyle=\tiny\color{polyagray},
  numbersep=8pt,
  frame=single,
  framerule=0.4pt,
  rulecolor=\color{polyagray},
  breaklines=true,
  showstringspaces=false,
  tabsize=4,
}
\lstset{style=polya-python}
\title{Root Finding \& Interpolation}
\author{uber-polya}
\date{2026-02-26}

\begin{document}

\maketitle
\tableofcontents
\newpage

% ── Phase A: Problem Formulation ─────────────────────────────────────
\section{Problem Formulation}

\begin{polyamodel}

\textbf{Problem Type}: Find \hfill
\textbf{Domain}: Numerical Methods


\end{polyamodel}

% ── Universe ──────────────────────────────────────────────────────────
\subsection{Universe}

\begin{itemize}
  \item $Root Finding: \{methods, best_method, best_root_degC, best_residual\}$
  \item $Interpolation: \{cubic_spline_max_error, linear_max_error, lagrange_max_error, max_spline_vs_linear, max_spline_vs_lagrange\}$
  \item $Integration: \{trapezoidal, simpson, gaussian_quad\}$
\end{itemize}

% ── Variables ─────────────────────────────────────────────────────────
\subsection{Variables}

\begin{longtable}{lll}
\toprule
\textbf{Symbol} & \textbf{Meaning} & \textbf{Type / Range} \\
\midrule
\endhead
$x$ & decision variable & $\mathbb{{R}}$ \\
\bottomrule
\end{longtable}

% ── Structure ─────────────────────────────────────────────────────────
\subsection{Structure}

Given 8 measured data points of temperature (0--700 degC) vs. thermal expansion coefficient alpha(T), which increases nonlinearly from \textasciitilde{}9 to \textasciitilde{}16 x10\textasciicircum{}-6 /degC:

1. **Root finding**: Find the temperature where alpha(T) = 12.0 x10\textasciicircum{}-6 /degC by solving g(T) = spline(T) - target = 0
2. **Method comparison**: Compare 4 root-finding methods -- bisection, Newton (finite-difference), secant, and Brent
3. **Interpolation comparison**: Fit the data using linear, cubic spline, and Lagrange interpolation
4. **Numerical integration**: Integrate alpha(T) over [100, 500] degC using trapezoidal, Simpson, and Gaussian quadrature

% ── Mapping ───────────────────────────────────────────────────────────
\subsection{Real-World Mapping}

\begin{longtable}{ll}
\toprule
\textbf{Real-World Concept} & \textbf{Mathematical Object} \\
\midrule
\endhead
problem input & $instance parameters$ \\
\bottomrule
\end{longtable}

% ── Constraints ───────────────────────────────────────────────────────
\subsection{Constraints}

\begin{enumerate}
  \item $Bisection**: Guaranteed convergence (linear rate), requires bracket$
  \hfill \textit{()}
  \item $Newton's method**: Quadratic convergence near root, needs derivative (approximated via finite difference)$
  \hfill \textit{()}
  \item $Secant method**: Superlinear convergence (~1.618), no derivative needed$
  \hfill \textit{()}
  \item $Brent's method**: Combines bisection reliability with superlinear speed$
  \hfill \textit{(the gold standard)}
  \item $Cubic spline**: C2-continuous interpolation that passes exactly through data points$
  \hfill \textit{()}
  \item $Gaussian quadrature**: Exact for polynomials up to degree 2n-1 with n points$
  \hfill \textit{()}
\end{enumerate}

% ── Objective / Claim ─────────────────────────────────────────────────
\subsection{Claim}


% ── Approach ──────────────────────────────────────────────────────────
\subsection{Approach}

\begin{description}
  \item[Suggested approach:] Cubic Spline + Brent/Bisection/Newton/Secant root finding
  \item[Complexity class:] 
  \item[Available tools:] Cubic Spline + Brent/Bisection/Newton/Secant root finding
\end{description}
% ── Phase B: Solution ────────────────────────────────────────────────
\section{Solution}

\begin{polyaresult}

\textbf{Answer}: Solution computed


\textbf{Optimal}: No / Unknown \hfill
\textbf{Feasible}: No
\textbf{Algorithm}: Cubic Spline + Brent/Bisection/Newton/Secant root finding \quad $$

\textbf{Time}: 0.011280923965387046s


\end{polyaresult}

% ── Solution Details ──────────────────────────────────────────────────
\subsection{Solution Details}

Solution computed successfully.

% ── Verification ──────────────────────────────────────────────────────
\subsection{Verification}

\begin{longtable}{lcl}
\toprule
\textbf{Check} & \textbf{Status} & \textbf{Detail} \\
\midrule
\endhead
All Methods Find Root & \passmark & Passed \\
Roots Agree & \passmark & Passed \\
Residual Small & \passmark & Passed \\
Root In Range & \passmark & Passed \\
Spline Passes Through Data & \passmark & Passed \\
Integration Methods Agree & \passmark & Passed \\
Integration Positive & \passmark & Passed \\
Brent Fewest Or Tied & \passmark & Passed \\
\bottomrule
\end{longtable}

% ── Phase C: Interpretation ──────────────────────────────────────────
\section{Interpretation}

% ── The Question & The Answer ─────────────────────────────────────────
\subsection{The Question}
Given 8 measured data points of temperature (0--700 degC) vs.

\subsection{The Answer}

\begin{polyainsight}[title={Bottom Line}]
A feasible solution was found.
\end{polyainsight}

% ── What This Means ───────────────────────────────────────────────────
\subsection{What This Means}

- Root (Brent): \textasciitilde{}244 degC where alpha(T) = 12.0 x10\textasciicircum{}-6 /degC
- Method agreement: All 4 methods agree to within 0.1 degC
- Brent iterations: Fewest (or tied) compared to bisection
- Residual: |g(root)| < 1e-8 for Brent's method
- Integration: All 3 quadrature methods agree within 1\%
- 8 verification checks: All pass

% ── Sensitivity Analysis ──────────────────────────────────────────────

% ── Figures ───────────────────────────────────────────────────────────

% ── Recommendations ───────────────────────────────────────────────────
\subsection{Recommendations}

\begin{enumerate}
  \item Review the model assumptions to ensure real-world fidelity.
\end{enumerate}

% ── Limitations ───────────────────────────────────────────────────────
\subsection{Limitations}

\begin{polyawarnbox}
\begin{itemize}
  \item Model assumes deterministic parameters; real-world variability not captured.
  \item Solution quality depends on the accuracy of input data.
\end{itemize}
\end{polyawarnbox}

% ── Appendix ─────────────────────────────────────────────────────────

\end{document}